%Motivation and research interests
How do simple local interactions give rise to collective behavior that no single constituent possesses? This question of \textit{emergent phenomena} (``More is Different'') is why I want a PhD in physics. I want to study strongly correlated quantum matter, especially quantum spin liquids, using ideas from quantum information, gauge theory, and topology. My long-term goal is a career in academic research, where I can pursue fundamental questions about quantum matter and teach and mentor students.

%Academic record
I completed my B.S. in physics at National Tsing Hua University in June 2024 with a graduation GPA of 4.1 on a 4.3 scale, and received the university's Academic Achievement Award three times. I then completed my M.S. there in June 2026 with a GPA of 4.25, and my thesis, \textit{Analysis of the Kitaev Model from the Perspective of Quantum Entanglement}, was graded A+. In my core condensed matter and theory courses, I earned A+ in Solid State Physics I and II and Statistical Mechanics I and II, A and A+ in Condensed Matter Physics I and II, and A+ in Quantum Mechanics I and II, Quantum Information, and Superconductivity. I ranked first in my class in several graduate courses, including Quantum Mechanics, Statistical Mechanics, and Electrodynamics. Before college I had already taught myself university-level mechanics, electromagnetism, relativity, quantum mechanics, and statistical mechanics, and I reached the final training camp of the Physics Olympiad selection.

%Research experience
I found my field when I attended Prof.~Yi-Ping Huang's talk on quantum spin liquids, phases with no local order parameter whose properties are encoded in long-range entanglement. I joined his group as an undergraduate. My first project was to express the ground state of the Kitaev honeycomb model as a tensor network. I taught myself frustrated magnetism and tensor networks for it, and the topic took shape through discussions with my advisor rather than being assigned. The project earned me an NSTC undergraduate research scholarship.

I stayed in the group for my master's degree and began collaborating with Dr.~Sungkit Yip at Academia Sinica, after Prof.~Huang introduced us during a 2024 summer internship. We first studied how topological order in the Kitaev ground state arises from many-body entanglement, working from a two-spin model to a chain and then to the honeycomb lattice. We then turned to the density matrix itself, and this became the core of my thesis and of a paper on which I am first author, published in \textit{Physical Review B}. We derived the exact density matrix of the Kitaev spin liquid directly in the physical spin Hilbert space, in terms of spin-string operators organized by equivalence classes that depend only on the strings' endpoints. Using this structure with the Gauss law of the emergent $\mathbb{Z}_2$ gauge theory and the model's exact 1-form Wilson symmetry, we showed that the reduced density matrix splits into $2^{|\partial A|/2-1}$ identical blocks. This solves the entanglement spectrum exactly and explains its degeneracy. The method also works for subsystems with an odd number of sites, where the usual Majorana construction is ill-defined. {\color{red} my own contribution} The work required lattice gauge theory, higher-form symmetries, and quantum-information tools, and I used symmetry-resolved exact diagonalization and the Lanczos method to check my results numerically. It also taught me to change direction when a question stops being productive and to state clearly what remains open. For example, we could not find a 1-form 't~Hooft symmetry operator independent of the Wilson loops, and we said so in the paper.

%Current and future research 
I have continued this line of work on quantum magic, measured by the stabilizer R\'enyi entropy, in the Kitaev model, starting with the one-dimensional chain and building a numerical library for it. I also want to learn which structures of quantum spin liquids survive at finite temperature or in mixed states, where the density matrix itself becomes the natural object. In graduate school I want to widen my toolkit beyond exactly solvable models to numerical methods for strongly interacting systems and to other correlated and topological phases.

{\color{red} Why: faculty in condensed matter theory I'm interested in, topic or paper, department strengths, centers, or collaborations that matter to you, Yuan-Ming Lu,  Nandini Trivedi,}

%Other qualities
I have worked to support other students as well as to study physics. I was a teaching assistant for Applied Mathematics, and since October 2022 I have tutored at my university's Counseling Center, supporting students with disabilities in their coursework. {\color{red} specific example} In September 2023 I received the university's Excellent Tutor Award. This work taught me to find where a student's understanding is actually breaking down and to explain the same idea in different ways. I also play the erhu, a traditional Chinese instrument, in my university's Chinese Orchestra, which won first place in the 2023 national student chamber-music competition. I expect these experiences to help me as a collaborator, a teaching assistant, and a member of a diverse and inclusive graduate community.